% ch05b —— §5.4–5.7：辛形式的构造、微分方程、度量及其大 $R$ 渐近、与物理度量的比较（原文 wcf-paper.tex 2545–3001 行）
\subsection{辛形式的构造}
\label{sec:symplectic-form}

至此，我们已经解出黎曼--希尔伯特问题，得到了一个映射 $\CX:\CM_u\to T_u$，它满足渐近条件 式 (5.3)、跳变条件 式 (5.7) 与实条件 式 (5.9)。% TODO 合稿改 \eqref：式 (5.3)(5.7)(5.9) 属前面各块
现在让 $u$ 变动，便得到一个映射 $\CX:\CM\to T$。于是我们拉回 $T$ 上典范的逐纤维辛形式，在 $\CM$ 上构造出一个复 $2$-形式 $\hsymp(\zeta)$：
\begin{equation} \label{eq:omega-pullback}
\hsymp(\zeta) = \frac{1}{4 \pi^2 R} \CX^* \hsymp^T = \frac{1}{8 \pi^2 R} \epsilon_{ij} \frac{d\CX_{\gamma^i}}{\CX_{\gamma^i}} \wedge \frac{d\CX_{\gamma^j}}{\CX_{\gamma^j}}.
\end{equation}
由式 \eqref{eq:omega-pullback} 可以直接得到 $\hsymp(\zeta)$ 的几条性质：
\begin{itemize}
\item 尽管 $\CX$ 关于 $\zeta$ 只是分段解析的，$\hsymp(\zeta)$ 却是真正解析的（因为诸不连续量 $S_\gamma$ 都是辛同胚，也就是说它们保持 $\hsymp^T$）。

\item 利用式 \eqref{eq:reality-X}，有 $\hsymp(- 1 / \bar\zeta) = \overline {\hsymp(\zeta)}$。

\item 当 $\zeta \to 0,\infty$ 时，利用 $\CX$ 的渐近行为 式 (5.3) 与 $\hsymp^\sf(\zeta)$ 的显式表达式 式 (3.15)，可以确定 $\hsymp(\zeta)$ 的行为。% TODO 合稿改 \eqref：式 (5.3)、(3.15) 属前面各块
我们发现，两种情形下 $\hsymp(\zeta)$ 都只有单极点，留数为
\begin{equation} \label{eq:omega-residues}
\Res_{\zeta = 0} \hsymp(\zeta) = \frac{i}{8 \pi} \Upsilon_0^*\inprod{dZ, d\theta}, \quad \Res_{\zeta = \infty} \hsymp(\zeta) = -\frac{i}{8 \pi} \Upsilon_\infty^*\inprod{d\bar Z, d\theta}.
\end{equation}

\item 利用 $\lim_{R \to \infty} \CX = \CX^\sf$ 可知，当 $R$ 充分大时 $\hsymp(\zeta)$ 非退化（在全纯意义下）。
\end{itemize}
这些性质在我们构造超凯勒度量时都将起重要作用。

\subsection{微分方程}
\label{sec:differential-problem}

我们的黎曼--希尔伯特问题，是用一些普适的不连续因子（它们局部地与规范理论的一切参数无关）来表述的，与之配套的渐近行为则由诸函数 $\CX^\sf_\gamma$ 给出，而这些函数对参数的依赖方式极为简单。本节将依照一套标准做法说明：这意味着解 $\CX$ 满足一族微分方程。

正如我们将看到的，这些方程的物理意义相当透明。其中一组表达了这样一件事实：解黎曼--希尔伯特问题所得的诸函数 $\CX(\zeta)$，在复结构 $J^\pz$ 下关于 $\CM$ 全纯。这组方程对于构造超凯勒度量是本质性的。另一对则描述重整化群流与一个 $U(1)_R$ 对称性的作用。它们对于把度量与 KS 穿墙公式联系起来很重要。

极为相似的一族方程，曾在「$tt^*$ 几何」的故事中起过关键作用；那一套东西出现在有质量的 $\CN=(2,2)$ 二维理论中 \cite{Cecotti:1993rm,Cecotti:1992qh,Cecotti:1991me,Cecotti:1993vy}。

我们首先回顾一点：我们的黎曼--希尔伯特问题在 $\IC\IP^1$ 上的解 $\CX$ 只是分段解析的，沿 $\IC\IP^1$ 中的各条射线 $\ell$，它带有 $\CX \to \CX S_\ell$ 形式的跳变。因此转而考虑\footnote{这是标准记号，但在我们的语境中它多少有些暗示意味，所以这里给出较长的说明。用 $\zeta \partial_\zeta$ 作用，使映射 $\CX$ 作无穷小变动，得到 $T$ 上的一个向量场，我们称之为 $\zeta \partial_\zeta \CX$。然后借助 $\CX$ 把它拉回，得到 $\CM_u$ 上的向量场 $\CA_\zeta$。我们把这一拉回运算记作 $\CX^{-1}$；由于我们采用的复合约定是非标准的，这个 $\CX^{-1}$ 出现在右边而不是左边，这使得我们的方程与黎曼--希尔伯特问题的通常形式一致；事实上，正是为了这一致性，我们在第 5 节才采用非标准的约定。% TODO 合稿改 \ref（Section \ref{sec:constructing-moduli}）
在局部坐标下，人们会写成
\begin{equation}
\CA_\zeta =  \frac{\partial \CX^i}{\partial \zeta}
\left[ \left(\frac{\partial \CX}{\partial \theta}\right)^{-1}
\right]^i_j \frac{\partial}{\partial \theta^j}.
\end{equation}
}
\begin{equation}
\CA_\zeta := \zeta \partial_\zeta \CX \CX^{-1}.
\end{equation}
$\CX$ 沿 BPS 射线的诸不连续量在 $\CA_\zeta$ 中相互抵消，因此 $\CA_\zeta$ 关于 $\zeta$ 是真正解析的——可能除了 $\zeta = 0, \infty$ 之外，那里 $\CX$ 变得奇异。于是我们可以把 $\CX$ 看作关于 $\zeta$ 的一个常微分方程的解，
\begin{equation} \label{eq:zeta-eq-general}
\zeta \partial_\zeta \CX = \CA_\zeta \CX.
\end{equation}
借助关于 $\CX$ 的渐近信息，可以把这一方程描述得相当具体。首先注意 $\CX^\sf$ 满足同一形式的方程。为写出它，先引入 $\CM_u$ 上的两个向量场，
\begin{equation}
\CA_\zeta^{(-1),\sf} := i \pi Z \cdot \partial_\theta, \quad \CA_\zeta^{(1),\sf} := i \pi \bar{Z} \cdot \partial_\theta.
\end{equation}
于是有
\begin{equation} \label{eq:zeta-eq-sf}
\zeta \partial_\zeta \CX^\sf = \CA^\sf_\zeta \CX^\sf,
\end{equation}
其中
\begin{equation}
\CA^\sf_\zeta = \frac{1}{\zeta} \CA_\zeta^{(-1),\sf} + \zeta \CA_\zeta^{(1),\sf}.
\end{equation}
要点在于，$\CA^\sf_\zeta$ 对 $\zeta$ 的依赖非常简单：只是在 $\zeta = 0,\infty$ 处各有一个单极点。我们可以把这一信息转换成关于 $\CA_\zeta$ 的信息，因为由式 (5.3) 知道 $\Upsilon = \CX (\CX^\sf)^{-1}$ 在 $\zeta = 0, \infty$ 两处都保持有限。% TODO 合稿改 \eqref：式 (5.3) 属前面各块
这说明 $\CA_\zeta$ 在 $\zeta = 0, \infty$ 处同样只有单极点，甚至还能定出留数，
\begin{equation}
\CA_\zeta = \frac{1}{\zeta}\CA_\zeta^{(-1)} + \CA_\zeta^{(0)} +  \zeta \CA_\zeta^{(1)},
\end{equation}
其中
\begin{equation}
\CA_\zeta^{(-1)} = \Upsilon_0 \CA_\zeta^{(-1),\sf} \Upsilon_0^{-1}, \quad \CA_\zeta^{(1)} = \Upsilon_\infty \CA_\zeta^{(1),\sf} \Upsilon_\infty^{-1}.
\end{equation}
因此我们看到，式 \eqref{eq:zeta-eq-general} 在 $\IC\IP^1$ 上定义了一个平坦联络 $\zeta \partial_\zeta - \CA_\zeta$，其取值在 $\CM_u$ 上向量场构成的无穷维代数中，并在 $\zeta = 0, \infty$ 处有秩 $1$ 的非正则奇点。$\CX$ 是这个联络的平坦截面。

于是，我们关于黎曼--希尔伯特问题的解，直接导致 $\IC\IP^1$ 上一个平坦联络的构造。事实上，这在常微分方程理论中是标准手法。我们得到的联络在 $\zeta = 0$ 与 $\zeta = \infty$ 处有非正则奇点，因而呈现斯托克斯现象。黎曼--希尔伯特构造的优点之一是：斯托克斯因子很容易确定——它们就是进入黎曼--希尔伯特问题的那些不连续量 $S_\ell$。

上面的讨论有一个重要的推广。我们面对的并非单个黎曼--希尔伯特问题，而是一整族这样的问题，它们随一些额外参数而变动。这些参数包括 $\vm$ 上的坐标 $u^i$，以及标度 $\Lambda$、$S^1$ 的半径 $R$，也许还有一些裸规范耦合 $\tau^0$。（眼下我们\emph{不}引入质量参数；但见下文第 6 节。）% TODO 合稿改 \ref（Section \ref{sec:masses}）
我们引入通用记号 $t^n$ 来统指这一切参数。

重要的是，定义黎曼--希尔伯特问题的诸不连续量 $S_\ell$ 不依赖于任何 $t^n$。因此，正如上面处理对 $\zeta$ 的依赖时那样，我们考虑
\begin{equation}
\CA_n := \partial_{t^n} \CX \CX^{-1}.
\end{equation}
与前面一样，$\CX$ 的诸不连续量相互抵消，所以 $\CA_n$ 关于 $\zeta$ 解析，只在 $\zeta = 0, \infty$ 处例外。同样与前面一样，我们可以先检查 $\CA^\sf_n := \partial_{t^n} \CX^\sf (\CX^\sf)^{-1}$ 的行为，以此控制奇点附近的行为。对我们的所有 $t^n$ 都有
\begin{equation}
\CA^\sf_n = \frac{1}{\zeta} \CA^{(-1),\sf}_n + \zeta \CA^{(1),\sf}_n
\end{equation}
其中 $\CA^{(\pm 1),\sf}_n$ 是些简单的向量场；然后和前面一样，利用 $\Upsilon$ 当 $\zeta \to 0,\infty$ 时有限这一事实，得到
\begin{equation}
\CA_n = \frac{1}{\zeta} \CA_n^{(-1)} + \CA_n^{(0)} + \zeta \CA_n^{(1)},
\end{equation}
其中
\begin{equation}
\CA_n^{(-1)} = \Upsilon_0 \CA_n^{(-1),\sf} \Upsilon_0^{-1}, \quad \CA_n^{(1)} = \Upsilon_\infty \CA_n^{(1),\sf} \Upsilon_\infty^{-1}.
\end{equation}
再把式 \eqref{eq:zeta-eq-general} 也算进来，我们得到的完整方程组是
\begin{align}
\partial_{u^j} \CX &= \left( \frac{1}{\zeta} {\CA^{(-1)}_{u^j}} + \CA^{(0)}_{u^j} \right) \CX, \label{eq:conn-holomorphy-1} \\
\partial_{\bar{u}^{\bar j}} \CX &= \left( {\CA}^{(0)}_{\bar{u}^{\bar j}} + \zeta {\CA}^{(1)}_{\bar{u}^{\bar j}} \right) \CX, \label{eq:conn-holomorphy-2} \\
\Lambda \partial_\Lambda \CX &= \left( \frac{1}{\zeta} {\CA_\Lambda^{(-1)}} + \CA_\Lambda^{(0)} \right) \CX, \label{eq:lambda-equation} \\
\bar\Lambda \partial_{\bar\Lambda} \CX &= \left( \CA_{\bar \Lambda}^{(0)} + \zeta \CA_{\bar \Lambda}^{(1)} \right) \CX, \label{eq:lambda-bar-equation} \\
R \partial_R \CX &= \left( \frac{1}{\zeta} {\CA_R^{(-1)}} + \CA_R^{(0)} + \zeta \CA_R^{(1)} \right) \CX, \label{eq:R-equation} \\
\zeta \partial_\zeta \CX &= \left( \frac{1}{\zeta} {\CA_\zeta^{(-1)}} + \CA_\zeta^{(0)} + \zeta \CA_\zeta^{(1)} \right) \CX. \label{eq:zeta-equation}
\end{align}
人们还会得到一些额外的关系
\begin{equation} \label{eq:M-relations}
\CA_R^{(-1)} = - \CA_\zeta^{(-1)}, \quad \CA_R^{(1)} = \CA_\zeta^{(1)},
\end{equation}
它们来自这样一件事实：$\CX^\sf$ 在 $\zeta \to 0$ 时被 $\zeta \partial_\zeta + R \partial_R$ 零化，在 $\zeta \to \infty$ 时被 $\zeta \partial_\zeta - R \partial_R$ 零化。

我们的方程到这里就构造完了。附录 D 讨论了在给定 $\CX$ 在 $\zeta = 0$ 附近的渐近展开时，如何把这些方程写得更具体。% TODO 合稿改 \ref（Appendix \ref{app:diff-asymptotics}）
本节最后给出几点评论：

\begin{itemize}
\item
既然辛形式 $\hsymp(\zeta)$ 是由 $\CX$ 构造出来的，式 \eqref{eq:zeta-equation} 与式 \eqref{eq:R-equation} 就平凡地蕴含了 $\hsymp(\zeta)$ 对 $\zeta$ 与 $R$ 的依赖所满足的方程，其形式为
\begin{align}
\zeta \partial_\zeta \hsymp &= \left( \frac{1}{\zeta} \CL_{\CA_\zeta^{(-1)}} +
\CL_{\CA_\zeta^{(0)}} + \zeta \CL_{\CA_\zeta^{(1)}} \right)  \hsymp,
\\ R \partial_R (R\hsymp) &= \left( \frac{1}{\zeta} \CL_{\CA_R^{(-1)}} +
\CL_{\CA_R^{(0)}} + \zeta \CL_{\CA_R^{(1)}} \right) (R \hsymp). \label{eq:omega-R}
\end{align}
回想式 (3.4) 中 $\hsymp(\zeta) = - \frac{i}{2 \zeta}\omega^+ + \omega_3 - \frac{i}{2} \zeta \omega_-$，把这些方程按 $\zeta$ 的幂展开，便可导出超凯勒形式 $\vec\omega$ 的一些有趣的微分方程。% TODO 合稿改 \eqref：式 (3.4) 属前面各块

\item
回想一下：黎曼--希尔伯特问题的解 $\CX$ 带有 $\CX \to b \CX$ 的任意性，其中 $b$ 是 $\CM_u$ 的任一微分同胚。这一任意性导致联络 $\CA$ 的、与 $\zeta$ 无关的规范变换。有几种特别方便的规范。其一是使 $\CA_R^{(0)} = 0$ 的规范。由式 \eqref{eq:omega-R} 可知，在这一规范下，$R \omega_3$ 到每个 $\CM_u$ 上的限制与 $R$ 无关（因而等于其 $R \to \infty$ 极限，即 $- \frac{1}{8 \pi^2} \inprod{d\theta, d\theta}$）。我们的积分方程 式 (5.11) 所选取的规范是否正是这一规范，将会很有意思。% TODO 合稿改 \eqref：式 (5.11) 属前面各块
如果我们允许 $b$ 是\emph{复化}的微分同胚，那么至少在形式上还可以选取使 $\Upsilon_0 = 1$ 的规范，从而 $\CA^{(-1)} = \CA^{(-1),\sf}$；这是 \cite{Cecotti:1993rm} 的「拓扑规范」的类比（对偶地，$\Upsilon_\infty = 1$ 将是「反拓扑规范」）。

\item
我们这些方程的两个线性组合有简单的物理意义：它们表达了理论在整体标度变换与 $R$ 对称变换下的不变性。为此首先注意
\begin{equation}
\left(a^I \partial_{a^I} + \Lambda \partial_\Lambda\right) Z_\gamma = Z_\gamma
\end{equation}
对一切 $\gamma \in \clat$ 成立。由此可得
\begin{align}
\left(R \partial_R - a^I \partial_{a^I} - \bar{a}^I \partial_{\bar{a}^I} - \Lambda \partial_\Lambda - \bar{\Lambda} \partial_{\bar{\Lambda}} \right) \CX^\sf &= 0, \\
\left(\zeta \partial_\zeta  + a^I \partial_{a^I} - \bar{a}^I \partial_{\bar{a}^I} + \Lambda \partial_\Lambda - \bar{\Lambda} \partial_{\bar{\Lambda}} \right) \CX^\sf &= 0.
\end{align}
这些方程可以解释为半平坦几何的（反常）标度不变性与 $R$ 对称不变性。它们在诸 $\CA^\sf_n$ 之间蕴含一些关系（只需作替换 $\partial \to \CA$ 即可），这些关系又给出诸 $\CA_n^{(\pm 1)}$ 之间的关系：我们发现
\begin{align}
\left(R \partial_R - a^I \partial_{a^I} - \bar{a}^I \partial_{\bar{a}^I} - \Lambda \partial_\Lambda - \bar{\Lambda} \partial_{\bar{\Lambda}} \right) \CX &= \Delta \CX, \label{eq:rescaling-1} \\
\left(\zeta \partial_\zeta  + a^I \partial_{a^I} - \bar{a}^I \partial_{\bar{a}^I} + \Lambda \partial_\Lambda - \bar{\Lambda} \partial_{\bar{\Lambda}} \right) \CX &= \Delta' \CX, \label{eq:rescaling-2}
\end{align}
其中 $\Delta$、$\Delta'$ 是 $\CM_u$ 上与 $\zeta$ 无关的向量场。通过规范变换可以令 $\Delta = 0$、$\Delta' = 0$。确实，我们的积分方程自动选取了合适的规范：附录 C 给出的 $R$ 很大时的递推解，逐项满足式 \eqref{eq:rescaling-1} 与式 \eqref{eq:rescaling-2}，且 $\Delta = \Delta' = 0$。% TODO 合稿改 \ref（Appendix \ref{app:integral-asymptotics}）
所以，在某种意义下，标度不变性与 $R$ 对称不变性在瞬子修正之后仍然存活。

\item
式 \eqref{eq:R-equation} 与式 \eqref{eq:zeta-equation} 之间的相容性，连同关系 \eqref{eq:M-relations}，蕴含四元组 $(\CA_\zeta^{(\pm 1)}, \CA_\zeta^{(0)}, \CA_R^{(0)})$ 对 $R$ 的依赖满足一组非线性微分方程。正如附录 D 所解释的，这些方程是 Nahm 方程的一个形变；% TODO 合稿改 \ref（Appendix \ref{app:diff-asymptotics}）
该四元组的大 $R$ 展开，可以直接通过迭代求解这些方程得到。它们是研究我们的构造在\emph{小} $R$ 处行为的一种可能工具。Nahm 方程的一般解会在某个有限的 $R$ 处变奇异，我们预期我们的问题也是如此。不过我们预期，这里所描述的这些特解——它们由 $\CN=2$、$d=4$ 场论中的 BPS 简并度所决定——实际上对一切 $R$ 都\emph{是}正则的。可能这会对 $\CN=2$ 理论可能的 BPS 谱与红外预势给出一个有趣的限制。\cite{Cecotti:1993rm} 中曾用极为相似的策略来限制 $d=2$ 理论的性质。

\item
本节的讨论为穿墙公式的作用提供了一个新的视角。方程组 \eqref{eq:conn-holomorphy-1}–\eqref{eq:zeta-equation} 描述了 $\IC\IP^1 \times \CP$ 上的一个平坦联络，其中 $\CP$ 是以诸 $t^n$ 为坐标的参数空间。这个平坦联络可以等价地看作 $\IC\IP^1$ 上联络的一个\emph{等单值形变族}，在 $\zeta=0, \infty$ 处有秩 $1$ 的非正则奇点。在每个 $t \in \CP$ 处，$\IC\IP^1$ 上该联络的斯托克斯数据由 KS 因子给出。利用沿 $\CP$ 的平行移动可以证明：非正则奇点处的斯托克斯数据在适当的意义下是「不变的」。确切地说：在 $\zeta$ 平面中任取一个凸扇形 $\CV$，则乘积
\begin{equation}
A_\CV = \prod_{\ell \subset \CV}^\ccwarrow S_\ell
\end{equation}
在使任何斯托克斯线 $\ell$ 都不进入或离开 $\CV$ 的任何 $t \in \CP$ 的变动下保持不变。把这一论断用到 $u$ 的变动上，就重新得到穿墙公式。
\end{itemize}

\subsection{度量的构造及其大 $R$ 渐近}
\label{sec:metric-construction}

至此我们已经构造了 $\CM$ 上的一族函数 $\CX_\gamma(\zeta)$，以及相应的全纯辛形式 $\hsymp(\zeta)$。正如第 3.2 节所讨论的，给定一个满足第 \ref{sec:symplectic-form} 节所列性质的 $\hsymp(\zeta)$，以及满足形如式 \eqref{eq:conn-holomorphy-1}、\eqref{eq:conn-holomorphy-2} 的「柯西--黎曼」方程的 $\CX_\gamma$，就存在 $\CM$ 上相应的超凯勒度量 $g$。% TODO 合稿改 \ref（第 3.2 节 = Section \ref{sec:twistor-general}）
这就是我们对 $g$ 的构造。

给定解出式 (5.13) 的精确函数 $\CX_\gamma(\zeta)$，我们可以把 $g$ 写成如下的封闭形式。% TODO 合稿改 \eqref：式 (5.13) 属前面各块
利用式 (5.13) 中核的展开（当 $|\zeta'/\zeta| < 1$ 时），得到 $\zeta \to 0$ 的渐近展开，% TODO 合稿改 \eqref：式 (5.13) 属前面各块
\begin{equation}
\log \CX_{\gamma} = \frac{1}{\zeta} F_{-1}^\gamma + F_0^\gamma +
\zeta F_1^\gamma + \CO(\zeta^2),
\end{equation}
其中
\begin{equation}
\begin{split}
F_{-1}^{\gamma} & = \pi R Z_\gamma, \\
F_0^\gamma & = i \theta_\gamma -  \frac{1}{4 \pi i} \sum_{\gamma'}
\Omega(\gamma';u) \langle \gamma,\gamma' \rangle
\int_{\ell_{\gamma'}} \frac{d\zeta'}{\zeta'}   \log (1-
\sigma(\gamma') \CX_{\gamma'}(\zeta')),\\
F_1^\gamma & = \pi R \bar Z_\gamma -  \frac{1}{2 \pi i}
\sum_{\gamma'} \Omega(\gamma';u) \langle \gamma,\gamma' \rangle
\int_{\ell_{\gamma'}} \frac{d\zeta'}{\zeta'^2}   \log (1-
\sigma(\gamma') \CX_{\gamma'}(\zeta')).
\end{split}
\end{equation}
然后代入 $\varpi$，即可提取出
\begin{equation}
\begin{split}
\omega_+ & = \frac{i}{2\pi^2 R} \epsilon_{ij} dF_{-1}^{\gamma^i}
\wedge d F_0^{\gamma^j}, \\
\omega_3 & = \frac{1}{8\pi^2 R} \epsilon_{ij} \left( 2
dF_1^{\gamma^i}\wedge d F_{-1}^{\gamma^j} + d F_0^{\gamma^i} \wedge
d
F_0^{\gamma^j} \right).\\
\end{split}
\end{equation}
由这些辛形式很容易得到 $g$。

现在来考虑 $g$ 在大 $R$ 处的行为。第 5.3 节讨论了诸 $\CX_\gamma$ 的大 $R$ 渐近，其中包括式 (5.14) 给出的首个 BPS 瞬子修正。% TODO 合稿改 \ref/\eqref：第 5.3 节 = Section \ref{sec:solving-rh}；式 (5.14) = \eqref{eq:X-first-approx}
现在把它翻译成对 $\hsymp(\zeta)$ 的修正。我们先算出 $\frac{d\CX_\gamma}{\CX_\gamma}$ 的修正：
\begin{equation} \label{eq:dX-correction}
\frac{d\CX_\gamma}{\CX_\gamma} = \frac{d\CX^\sf_\gamma}{\CX^\sf_\gamma} + \CI_\gamma + \cdots,
\end{equation}
其中
\begin{equation} \label{eq:correction-integral}
\CI_\gamma =  \frac{1}{4\pi i} \sum_{\gamma'} \degen(\gamma';u)
\langle \gamma,\gamma' \rangle \int_{\ell_{\gamma'}}
\frac{d\zeta'}{\zeta'} \frac{\zeta' + \zeta}{\zeta' - \zeta}
\frac{d\CX^\sf_{\gamma'}(\zeta')}{\CX^\sf_{\gamma'}(\zeta')}
\frac{\sigma(\gamma') \CX^\sf_{\gamma'}(\zeta')}{1- \sigma(\gamma')
\CX^\sf_{\gamma'}(\zeta')}.
\end{equation}
注意，正如所承诺的，$\CI_\gamma$ 当 $R \to \infty$ 时指数压低，因为在 $\ell$ 上，当 $R \to \infty$ 时 $\CX_{\gamma'}^\sf$ 指数地趋于零。式 \eqref{eq:dX-correction} 中的省略号表示多瞬子修正，它们被压得更低。因此 $\hsymp(\zeta)$ 的首个修正来自 $\frac{d\CX^\sf_\gamma}{\CX^\sf_\gamma}$ 与 $\CI_\gamma$ 之间的楔积。

为更显式地描述这一修正，方便的做法是逐个处理每个 $\gamma'$，并取一个辛基 $\{\gamma^1, \dots, \gamma^{2r}\}$，使 $\gamma' = q_{\gamma'} \gamma^1$。于是式 \eqref{eq:correction-integral} 中的积分与积分 式 (4.38) 本质上相同——后者在第 4.3 节给出了 $\CX_\magn$ 的瞬子修正。% TODO 合稿改 \eqref/\ref：式 (4.38) = \eqref{eq:ipm}；第 4.3 节 = Section \ref{sec:solving-xmagn}
像那里一样算出 $\hsymp(\zeta)$ 相应的修正，便得到
\begin{equation}
\hsymp(\zeta) = \hsymp^\sf(\zeta) + \sum_{\gamma' \in \clat} \hsymp^\inst_{\gamma'}(\zeta) + \cdots,
\end{equation}
其中
\begin{equation} \label{eq:Omega-inst-general}
\hsymp^\inst_{\gamma'}(\zeta) = - \degen(\gamma';u) \frac{1}{4 \pi^2 R}\frac{d \CX^\sf_{\gamma'}(\zeta)}{\CX^\sf_{\gamma'}(\zeta)} \left[ A^\inst_{\gamma'} +
\frac{1}{2} V^\inst_{\gamma'} \left(\frac{1}{\zeta} da_{\gamma'} - \zeta d\bar{a}_{\gamma'}\right)\right],
\end{equation}
而（对照式 (4.7)、式 (4.10)）% TODO 合稿改 \eqref：式 (4.7) = \eqref{eq:V-inst}，式 (4.10) = \eqref{eq:A-inst}
\begin{align} \label{eq:A-inst-general}
V^\inst_{\gamma'} &= \frac{R q_{\gamma'}^2}{2\pi} \sum_{n>0}
\sigma(n \gamma') e^{i n \theta_{\gamma'}} K_0(2 \pi R \abs{n
Z_{\gamma'}}), \\
A^\inst_{\gamma'} &= - \frac{R q_{\gamma'}^2}{4 \pi} \left( \frac{dZ_{\gamma'}}{Z_{\gamma'}} - \frac{d \bar Z_{\gamma'}}{\bar Z_{\gamma'}} \right) \sum_{n>0} \sigma(n \gamma') e^{i n \theta_{\gamma'}} \abs{Z_{\gamma'}} K_1(2 \pi R \abs{n Z_{\gamma'}}).
\end{align}
从这里出发，可以按 $\zeta$ 展开，以提取 $\omega_+$、$\omega_3$ 的首阶修正，从而得到 $g$ 的首阶修正。

\subsection{与物理度量的比较}
\label{sec:comparison}

我们已经在 $R$ 充分大时于 $\CM$ 上构造出超凯勒度量 $g$，现在把它的若干性质总结如下：

\begin{enumerate}
\item $g$ 连续；
\item 若一切 BPS 粒子都满足 $|Z| \to \infty$，则 $g$ 趋于半平坦度量 $g^\sf$；
\item $g$ 光滑，只有那些物理上预期的奇点例外——它们位于 $\vm$ 中诸奇异轨迹的上方；
\item $g$ 满足 $\vol(\CM_u) = \left(\frac{1}{R}\right)^{r}$；
\item 在复结构 $J_3$ 下，$(\CM, g)$ 可以等同于处于标准复结构的 SW 环面纤维化；在这一等同之下，全纯辛形式为 $\omega_+ = - \frac{1}{4 \pi} \inprod{dZ, d\theta}$。
\end{enumerate}

所有这些性质都与 $\IR^3 \times S^1$ 上物理度量的预期相符，后者在 \cite{Seiberg:1996nz} 中有描述。因此，最简单的自洽图像就是：我们构造的度量确实就是物理度量。（在秩 $1$ 情形，\cite{Seiberg:1996nz} 曾提出，借助丘成桐（Yau）定理的一个非紧类比，这些性质确实\emph{决定}了度量。更一般地，存在类似定理的可能性是相当大的。）

本节的余下部分将从我们的构造出发论证这些性质：
\begin{enumerate}
\item $g$ 的连续性由穿墙公式推出，这一点我们已经解释过。

\item 只需一看修正项的形式 \eqref{eq:Omega-inst-general} 即可：它们全都被 $R \abs{Z_{\gamma'}}$ 指数压低，因而当一切 $\abs{Z_{\gamma'}} \to \infty$ 时指数快地消失。

\item 在任何使 $R \abs{Z_\gamma} \to \infty$（对一切 $\gamma$）的极限下，瞬子贡献都指数压低，$g$ 趋于 $g^\sf$。这足以确立 $g$ 在大 $R$ 处的光滑性，\emph{除了}在某些奇异轨迹附近——那里有若干电荷为 $\gamma_i$ 的 BPS 粒子变成无质量（$Z_{\gamma_i} = 0$ 且 $\degen(\gamma_i;u) \neq 0$）。为理解这些点附近的行为，我们考虑一个标度极限：$R \to \infty$ 而保持 $R Z_{\gamma_i}$ 有限。在这一极限下，可以用一个只保留 BPS 射线 $\ell_{\gamma_i}$、丢掉其余一切射线的黎曼--希尔伯特问题来近似原来的问题。确实，其余一切不连续量都含有 $(1-\sigma(\gamma')\CX_{\gamma'})$ 形式的因子，而在这个标度极限下它们指数地趋近于 $1$。

在最简单的情形，即只有一个 $Z_\gamma = 0$ 时，我们总可以选取一个对偶标架，使 $\gamma$ 为电电荷。通过把某些角 $\theta$ 平移 $\pi$，还可以安排细化 $\sigma$ 在该标架下取标准形式 $\sigma = (-1)^{\gamma_e \cdot \gamma_m}$。于是我们就处在第 4 节所研究的情形中，那里找到了一个超凯勒度量，它除了一个周期性的 $q$ 个 $A_{q-1}$ 奇点阵列之外是光滑的。% TODO 合稿改 \ref（第 4 节 = Section \ref{sec:periodicnut}）
这与 $d=3$ 有效场论的预期一致：奇点出现在电荷为 $q$ 的超多重态 Kaluza--Klein 塔中某一个变成无质量的地方。

除了我们考察过的这些物理奇点——即一组相互局域的 BPS 粒子变成无质量之处——之外，还可能有超共形点，那里相互非局域的粒子同时变成无质量 \cite{Argyres:1995jj,Argyres:1995xn}。我们尚未分析这些奇点；尽管我们预期它们会很有意思，并且预期二次加细在其分析中将起重要作用。

\item 由于 $\CM_u$ 关于 $J_3$ 是一个复环面，其体积就是 $\frac{1}{r!} \int_{\CM_u} \omega_3^r$。另一方面，利用式 (3.4)，并注意到由式 \eqref{eq:omega-residues} 知 $\omega_{\pm}$ 在 $\CM_u$ 上的限制为零，便有% TODO 合稿改 \eqref：式 (3.4) 属前面各块
\begin{equation}
\vol(\CM_u) = \frac{1}{r!} \int_{\CM_u} \hsymp^r(\zeta) = \frac{1}{(4 \pi^2 R)^r r!} \int_{\CX(\CM_u)} (\hsymp^T)^r = \left(\frac{1}{R}\right)^r,
\end{equation}
正如所愿。

\item 复结构 $J_3$ 可以由 $\omega_1$ 与 $\omega_2$ 定出，只需 $J_3 = \omega_1^{-1} \omega_2$。但这一信息又由 $\hsymp(\zeta)$ 在 $\zeta = 0$ 处的留数给出；回想式 (3.4)，$\omega_+ = \omega_1 + i \omega_2$ 正是由% TODO 合稿改 \eqref：式 (3.4) 属前面各块
\begin{equation}
\omega_+ = 2i\,\Res_{\zeta = 0} \, \hsymp(\zeta).
\end{equation}
给出的。我们关于 $\hsymp(\zeta)$ 的渐近条件 式 (5.3) 恰好保证它与 $\hsymp^\sf(\zeta)$ 的留数相关：确实，我们正好有% TODO 合稿改 \eqref：式 (5.3) 属前面各块
\begin{equation}
\omega_+ = \Upsilon_0^* \omega_+^\sf.
\end{equation}
由此可知，只需作用以逐纤维微分同胚 $\Upsilon_0$，就能把 $(\CM, J_3)$ 等同于 $(\CM, J_3^\sf)$。正如第 3.3 节所解释的，$\CM$ 上的复结构 $J_3^\sf$ 正是 SW 环面纤维化的复结构。% TODO 合稿改 \ref（第 3.3 节 = Section \ref{sec:semiflat-twistor}）
此外，在这一等同之下，$\omega_+$ 被等同于式 (3.16) 给出的 $\omega_+^\sf$。% TODO 合稿改 \eqref：式 (3.16) 属前面各块
\end{enumerate}
